\subsection{Exact additive-coordinate lattices from one square-zero lift}
\label{sec:geometry-witt-additive-lattice}

This argument supplies the finite flat group calculations underlying
(RC1)--(RC2). It does not require a separate full-order tightness theorem.
We give the calculation over the valued test fields first, and retain
the chart-to-test-model identification in the application.

Let \(C\) be an algebraically closed complete rank-one field of mixed
characteristic \((0,p)\), \(R=\mathcal O_C\), \(v(p)=1\), and \(p>2\).
Choose \(w\) with \(v(w)=1/(p-1)\).
For a \(p\)-divisible group \(P/R\) let
\[
 T_P=R\otimes_{\mathbf Z_p}T_pP,\qquad
 M_P=\mathcal O(P[p])\subset R^{P[p](C)}.
\]
The inclusion uses evaluation on the specified generic points.
Write \(A_P\subset T_P^\vee\) for the coefficient lattice of those
linear functions on \(P[p](C)\) modulo \(pR\) that lift to \(M_P\).
Equivalently, after choosing a Tate basis, the row \(a\) belongs to
\(A_P\) when the function \(\sum_j a_j x_j\), with integer
representatives \(x_j\) of the \(\mathbf F_p\)-coordinates, belongs
to \(M_P+pR^{P[p](C)}\).
Let \(H_P=A_P^\vee\subset T_P\otimes_R C\).
The definition is intrinsic, since changes of Tate basis are linear
modulo \(p\). Denote the integral character-differential map by
\[
 a_P:T_{P^D}\longrightarrow\omega_P.
\]
A compatible choice of primitive roots identifies the natural Tate
duality \(T_{P^D}=T_P^\vee\). This identification is used only
with its actual Cartier pairing.

We use the annihilator and determinant arguments of
Section~\ref{sec:geometry-torsion-character-flags}, including the
case \(\mathbb F=\mathbf F_p\). In particular, for \(J=P[p]\),
\(h=\operatorname{ht} P\), and \(d=\dim P\),
\begin{equation}\label{eq:aal-bound}
 wT_P^\vee\subset A_P\subset T_P^\vee,\qquad
 (p-1)v(\det A_P)\ge d.
\end{equation}
The second statement follows from
\(2v(\det M_P)=p^h d\) and the evaluation determinant of
monomials in Smith-normalized additive functions.
Equality in the second statement gives generation of \(M_P\)
as an \(R\)-algebra by functions additive on the generic points
modulo \(pR\). The determinant and discriminant argument here
requires no principality of a polarization.

\paragraph{A division identity in length-two Witt vectors.}
Write \(R^\flat\) for the tilt, \(A_{\inf}=W(R^\flat)\),
and choose \(\underline p\in R^\flat\) with
\(\underline p^\sharp=p\). Put
\[
 \xi=[\underline p]-p,\quad S=A_{\inf}/(\xi^2),\quad
 V=A_{\inf}/(p^2,[\underline p]).
\]
We use the standard perfectoid identities
\(\ker\theta=(\xi)\) and \(R^\flat/(\underline p)\simeq R/pR\)
(Bhatt--Morrow--Scholze, \emph{Integral \(p\)-adic Hodge theory},
Publ.\ Math.\ IHES 128 (2018), \S3, especially Lemmas 3.9--3.10).
The ring \(S\) is a square-zero extension of \(R\) by \(\xi R\).
The element \(\xi\) maps to \(-p\) in \(V\), giving \(S\to V\).
Moreover \(V/pV=R/pR\) and multiplication by \(p\) induces
an isomorphism \(V/pV\to pV\). Indeed
\([\underline p]\) reduces to a nonzero element of the domain
\(R^\flat\); if \(px=[\underline p]y\bmod p^2\) in
\(A_{\inf}\), reduction modulo \(p\) gives \(y\in pA_{\inf}\),
and cancellation of \(p\) gives the claimed kernel.

Suppose \(F(X_1,\ldots,X_e)\in V[X_1,\ldots,X_e]\) has all
partial derivatives zero modulo \(p\). There is a polynomial
\(D_F\in (R/pR)[X_1,\ldots,X_e]\) such that, for any \(x\in V^e\)
with \(F(x)\in pV\),
\begin{equation}\label{eq:aal-witt-division}
 F(x)/p\bmod p=D_F(\bar x).
\end{equation}
Here division uses the preceding isomorphism, not an inverse
of \(p\) in \(V\). To prove this, write
\(F=\sum_m[a_m]X^{pm}+pQ(X)\).
Indeed any monomial with an exponent not divisible by \(p\)
has coefficient zero modulo \(p\), and every coefficient
modulo \(p\) has a Teichm\"uller representative. Choose
\(x_j=[b_j]+pu_j\) and set
\(z_m=a_m^{1/p}\prod_j b_j^{m_j}\) in the perfect ring \(R^\flat\).
For the integral polynomial
\[
 {\cal C}_p(Z)=\frac{\sum_m Z_m^p-(\sum_m Z_m)^p}{p}
\]
one has in \(V\)
\[
 \sum_m[z_m^p]=[(\sum_m z_m)^p]+p{\cal C}_p([z_m]).
\]
This follows by expanding \((\sum_m[z_m])^p\); the latter is
\([(\sum_m z_m)^p]\) modulo \(p^2\). Also \(x_j^p=[b_j^p]\).
The hypothesis \(F(x)\in pV\) says that
\((\sum_m z_m)^p\in(\underline p)\); its Teichm\"uller lift
is therefore zero in \(V\).
It follows that \(D_F=\bar Q+{\cal C}_p(
\overline{[a_m^{1/p}]}X^m)\). This identity is valid over
the nonreduced ring \(R/pR\).

\begin{theorem}[The exact additive lattice for good abelian reduction]
\label{thm:geometry-good-abelian-additive}
If \(P=B[p^\infty]\) for an abelian scheme \(B/R\), then
\[
 (p-1)v(\det A_P)=\dim B,\qquad
 A_{P^D}=wH_P,\qquad a_P(H_{P^D})=\omega_P.             \tag{AL1}
\]
The middle equality is in the naturally paired Tate spaces.
In particular \(M_P\) is generated by affine-modulo-\(p\) functions.
These conclusions concern the actual evaluation order and actual
Cartier character differentials, independently of an auxiliary
polarization or a chosen lift of \(B\).
\end{theorem}
\begin{proof}
First assume \(B\) principally polarized, of dimension \(g\).
The smooth moduli scheme of principally polarized abelian schemes
with full prime-to-\(p\) symplectic level at least three supplies a
principally polarized lift \(\widetilde B/S\). We use here the
classical smoothness theorem for Siegel moduli (Faltings--Chai,
\emph{Degeneration of Abelian Varieties}, Ch.~I).
The prime-to-\(p\) level exists over the strictly henselian \(R\)
and its root of unity lifts through the nilpotent ideal.
Formal smoothness gives a lift through \(S\to R\), since the base is
affine and local. No lift of an additional endomorphism algebra is
required.

Use the polarization to identify \(P=P^D\), and let
\(E\) denote the unimodular alternating matrix of the actual
Tate pairing. If \(t=(t_n)_n\) belongs to \(T_pP\), lift
\(t_n\) to \(\widetilde t_n\in\widetilde B(S)\).
The infinitesimal point
\([p^n]\widetilde t_n\) identifies with \(\xi b_n(t)\)
for \(b_n(t)\in\operatorname{Lie}B\).
Changing the lift changes \(b_n(t)\) by a multiple of \(p^n\);
multiplying a lift at level \(n+1\) by \(p\) proves compatibility
modulo \(p^n\). Addition of lifts proves additivity. Completeness
and separatedness of \(R\) give a map
\(b:T_P\to\operatorname{Lie}B\).
Write \(a:T_P\to\omega_B\) for the character-differential map
after the given polarization.

In integral dual bases for \(\omega_B\) and \(\operatorname{Lie}B\),
all rows of \(a\) belong to \(A_P\). Indeed differentiating the
universal pairing on \(B[p]\times B[p]\) at zero in its second
variable gives a differential over the first-variable coordinate
algebra, valued in \(\omega_B/p\omega_B\).
Its evaluation on a generic point is the specified character
differential modulo \(p\).

All rows of \(b\) belong to \(A_P\) as well. For precision, the
finite free algebra of \(B[p]\) decomposes into its components
over the closed residue points: use strict henselianity and lift
the special-fibre idempotents. Every such component has a
generic point, which extends to a section \(y\in B[p](R)\) by
properness. Choose a lift of this section to \(\widetilde B(S)\).
On the corresponding component, the universal point modulo \(p\),
translated by \(-y\), has formal coordinates at the identity.
These coordinates are nilpotent in the component algebra
modulo \(p\). In fact \(R/pR\) has only its closed prime, and
the finite component algebra has only the chosen closed prime;
its translated coordinates vanish there. Thus a single finite
nilpotence bound works for all its generic evaluations modulo
\(p\). For any chosen lifts of those points to \(\widetilde B(V)\),
the corresponding translated coordinates are nilpotent with at
most twice this bound, since \(p^2=0\) in \(V\).
Smooth formal coordinates at the identity are available over
the square-zero base; their evaluation on nilpotent tuples
uses only finitely many coefficients.

In these coordinates the map
\(X\mapsto [p]\widetilde y+[p]X\) is a tuple of formal series.
Its reduction modulo \(p\) has zero partial derivatives:
the differential of multiplication by \(p\) on a smooth
commutative group is zero in characteristic \(p\), at every
point by translation. Truncate beyond the uniform nilpotence
bound, so that the assertion and the evaluations hold for
polynomials. Their evaluations at the lifts of the torsion
points are the coordinates of \(\xi b_1(t)\), hence equal to
\(-p b_1(t)\) in \(V\). Equation
\eqref{eq:aal-witt-division} represents \(b_1(t)\bmod p\) by
polynomials in the universal translated coordinates modulo
\(p\) on this component. The component idempotents assemble
these into functions on \(B[p]\bmod p\); they lift to
\(\mathcal O(B[p])\) because this is an affine finite free
scheme. This proves row membership for \(b\). It does not
assert that the division by \(p\) of an arbitrary function
on a torsion scheme is regular.

For the determinant, there is an exact relation between the
two maps \(a,b\). Let \(\varepsilon=(1,\zeta_p,\zeta_{p^2},
\ldots)\) in \(R^\flat\) be the chosen cyclotomic sequence and
write \(\mu=[\varepsilon]-1\).
The other standard generator of \(\ker\theta\) is
\(\xi_{\rm cyc}=\mu/([\varepsilon^{1/p}]-1)\).
It follows that \(c=\theta(\mu/\xi)\) has
\(v(c)=1/(p-1)\).
We check the pairing relation using the rigidified Poincar\'e
biextension of \(\widetilde B\) and its lifted polarization;
this is a first-order calculation, not a comparison of
archimedean periods.

For \(N=p^n\), work over \(S/(N)\). Its ideal \((\xi)\) is
square-zero and is identified with \(R/NR\), since \(R\)
is \(p\)-torsion-free and \(S\) is the extension above.
Choose lifts \(x,y\) of \(t_n,s_n\) and put
\(u=Nx=\xi b_n(t)\), \(v'=Ny=\xi b_n(s)\).
Take a generator of the invertible module at \((x,y)\).
Its two \(N\)-fold biextension transports give generators at
\((u,y)\) and \((x,v')\). Relative to generators normalized to
reduce to \(1\) by the rigidifications along the axes, call
their ratios \(r_1,r_2\).
The quotient \(r_1/r_2\) reduces to the Weil pairing of
\(t_n,s_n\) (choose this convention; using its inverse changes
the common sign below). Transporting once more gives
\[
 r_1^N U_1=r_2^N U_2
\]
by the interchange axiom of the biextension, where \(U_1,U_2\)
are the transports along \(y\), respectively \(x\), of the
normalized generators. The fibre at \((u,v')\) is canonically
trivial to first order: the two axis rigidifications agree
there, since every term vanishing along both axes is a product
of two square-zero coordinates. The following evaluation of
\(U_1,U_2\) is useful. Transport of a normalized generator at
\((u,y)\) along the first variable is \(1\), because \(Nu=0\)
and the multiplication cocycle at two infinitesimal first
coordinates is zero by the axis rigidifications. Transport
along the second variable is identically \(1\) at \(u=0\);
its first derivative in \(u\) consequently depends only on
the reduction of \(y\). By the defining two transports for
the Weil pairing it is the derivative of
\(e_N(u,s_n)^{-1}=e_N(s_n,u)\).
The same argument at \((x,v')\), exchanging the variables,
gives the derivative of \(e_N(t_n,v')\). Thus, with the
convention \(a(t)=d\log e_N(t_n,-)\) modulo \(N\),
\[
 U_1=1+\xi\,a(s)(b(t)),\qquad
 U_2=1+\xi\,a(t)(b(s)).
\]
All first-order identities are independent of the choice of
normalized generators, whose \(N\)-th powers are unchanged
when multiplied by \(1+\xi z\).

On the other hand \(r_1/r_2\) differs from the lift
\([\varepsilon^{1/N}]^{E(t,s)\bmod N}\) by \(1+\xi z\).
Its \(N\)-th power is therefore
\(1+\xi c\,E(t,s)\) in \(S/(N)\).
Taking coefficients of \(\xi\) and passing through all \(n\)
gives the matrix identity
\[
 {\cal P}^{\,t}J{\cal P}=cE,\qquad
 {\cal P}=\begin{pmatrix}a\\ b\end{pmatrix},\qquad
 J=\begin{pmatrix}0&1_g\\-1_g&0\end{pmatrix}.            \tag{AL2}
\]
A reversed choice of the pairing convention replaces \(c\)
by \(-c\), which has the same valuation and generates the
same ideal. This calculation uses only the actual lifted
biextension, and recovers the source's alternating two-map
identity without an unspecified normalization factor.

Taking determinants gives \(v(\det{\cal P})=g/(p-1)\).
The row lattice of \({\cal P}\) is contained in \(A_P\);
\eqref{eq:aal-bound} forces equality. In particular the
monomial generation criterion applies. Also
\(H_P={\cal P}^{-1}R^{2g}\), and the first projection gives
\(a(H_P)=\omega_B\). From (AL2), \(wEH_P=A_P\)
in column notation for the dual lattice, since \(c/w\)
is a unit. This is exactly \(A_{P^D}=wH_P\) under the
canonical pairing and the principal polarization.

For any abelian scheme, apply Zarhin's classical trick to
\(B^4\times(B^\vee)^4\). The resulting principal polarization
is used only to prove (AL1) for that product. All three
identities descend to each factor: additive-coordinate
lattices of products are direct sums, as restriction to the
two factors proves one inclusion and addition of functions
proves the other. The definition of the differential map
and canonical Tate duality also commute with products and
isomorphisms. This proves (AL1) on \(B\) without transporting
its lattice by a nonintegral isogeny.
\end{proof}

\begin{corollary}[Induced groups and all coordinate quotients]
\label{cor:geometry-induced-additive}
At a height-one valued test point in
Lemma~\ref{lem:corrected-ht-test-tower}, the induced group \(P\)
satisfies (AL1) for the canonical Cartier pairing. The extended
original Weil pairing is an isomorphism \(P\simeq P^D\).
The actual function algebra of every coefficient-field
factor of \(P[p]\), and of every quotient of that factor by
a schematic coordinate subgroup, is generated by
affine-modulo-\(p\) functions.
\end{corollary}
\begin{proof}
Use the Raynaud extension of the given semistable degeneration
over the complete height-one DVR, after a finite extension
splitting the special torus, then base change to \(R\).
This is the same extension as in the complete-DVR
uniformization supplier: its formal completion is that of
the given semiabelian family (Faltings--Chai, Ch.~II;
Lan, \emph{Arithmetic compactifications of PEL-type Shimura
varieties}, the degeneration and uniformization theorems).
Its finite \(p^n\)-torsion identifies with the finite part
\(\mathcal H_n\) in that lemma: both finite schemes have the
same completion, and the finite part is recovered over the
complete DVR. Thus the compatible finite parts form an
extension of the good abelian \(p\)-divisible group by a
multiplicative \(p\)-divisible group.
The quotient of \(P\) by these finite parts is constant of
height the torus rank (the ranks at every level are those
of the finite-part and full generic torsions). A splitting
of the generic exact sequence of Tate modules extends to
a section of the constant quotient: each generic point of
a finite flat group extends over the integrally closed
valuation ring, and the group and compatibility identities
extend by generic density. Similarly a generic retraction
onto the multiplicative subgroup of the finite-part group
extends, by applying this argument to its Cartier dual,
whose quotient is constant. Hence, over \(R\), \(P\)
is a product of the good abelian group and multiplicative
and constant factors. This is an integral splitting over the
test ring only, and need not respect the coefficient action.

For \(\mu_{p^\infty}\) and \(\mathbf Q_p/\mathbf Z_p\) the
additive lattices are \(wR\) and \(R\), respectively, by
\eqref{eq:aal-bound} and the explicit characters; the
character-differential maps and duality in (AL1) hold directly.
Theorem~\ref{thm:geometry-good-abelian-additive} and the
product argument prove (AL1) for \(P\).
The morphism \(\lambda:P\to P^D\) constructed from the
original pairing is a generic isomorphism and both groups
have height \(2g\) and dimension \(g\).
On \(p\)-torsion the injection of evaluation orders given by
\(\lambda^*\) has equal determinant valuations on the two
sides, by the discriminant formula of the preceding section.
It is an equality of finite free lattices. Thus \(\lambda[p]\)
is an isomorphism, and the exact sequences of the successive
\(p\)-power kernels give the same assertion at every level.

For generation on coordinate quotients, embed the product
\(P[p]\) in the \(p\)-torsion of an abelian scheme, adding
ordinary elliptic schemes to supply the multiplicative and
constant factors. Such a good ordinary elliptic scheme exists
over \(R\) by lifting an ordinary elliptic curve over the
algebraically closed residue field; its \(p\)-divisible group
splits into these factors over \(R\) by the constant-section
argument. A coefficient-field factor and all its schematic
coordinate subgroups are finite flat closed subgroups.
For \(K\) one such subgroup, quotient the ambient abelian
scheme by \(K\); the quotient is an abelian scheme by the
finite flat subgroup quotient theorem. The specified factor
modulo \(K\) is a closed subgroup of its \(p\)-torsion:
the induced map is a proper monomorphism, by fppf descent
of the kernel. Restriction of functions is therefore
surjective, and it takes affine-modulo-\(p\) generators to
affine-modulo-\(p\) generators. Apply the theorem to that
ambient quotient. This proves generation for the original
coordinate subquotient, with no splitting of its flag.
\end{proof}

\paragraph{The literal image and kernel comparison.}
At a new correction prime, \(p>2\) is unramified in the
coefficient field. The original labelled modules are the
components of the split unramified coefficient algebra after
extension to \(R\). The fundamental Teichm\"uller characters
reduce modulo \(p\) to its labelled embeddings (the residue
label uniquely determines an embedding of the unramified
ring). The change from an \(\mathbf F_p\)-basis to these
characters has the invertible Moore matrix described in the
preceding section. Thus the \(H_\tau\) of (RC1) are exactly
the labelled components of \(H_P\), when \(M\) is the
canonical induced evaluation order on the corresponding
finite field factor. Use the original trace-pairing
identification between conjugate labels. Its different
is a unit at this prime.

Let \(q_*:T_\tau\otimes C\to O_{\tau,\eta}\) be the actual
Tate character-differential map and let \(i_*:L_{\tau,\eta}
\to T_\tau\otimes C\) be the adjoint of the conjugate
character map under the actual pairing. On the canonical
test models, \(q_*(T_\tau)=O_\tau^{\rm HT}\) and
\(i_*^{-1}(T_\tau)=W_\tau\). These are the image lattice
and its conjugate dual in (HF1)--(HF2); the latter equality
follows directly from dualizing the image map and the
unimodular pairing. Equations (AL1), with the extended
original polarization, imply
\[
 q_*(H_\tau)=O_\tau,\qquad
 i_*^{-1}(H_\tau)=w^{-1}L_\tau.                         \tag{AL3}
\]
For the second equality, \(w\) times the original pairing
identifies \(H_\tau\) with the dual of the conjugate \(H\);
testing the adjoint against the conjugate \(H\) is
equivalent to testing \(w\) times a tangent vector against
the entire old cotangent lattice, by the first equality.

Now use the actual finite-level lifts \(i,q\) of the maps
in (HF1), with \(m=R_p-2\ge2\).
By (HF2) and its adjoint, their differences from the actual
Tate maps satisfy
\[
 (i-i_*)W_\tau\subset p^m T_\tau,\qquad
 (q-q_*)T_\tau\subset p^m O_\tau^{\rm HT}.
\]
The graph definitions of \texttt{Gamdel} are
\[
 \Gamma_\tau=\{(z,t):wz\in W_\tau,\ t-iz\in T_\tau\},
 \qquad
 \Delta_\tau=\{(t,o):wt\in T_\tau,\
                    o-qt\in O_\tau^{\rm HT}\}.
\]
Since \(p^mR\subset wR\), the differences in the two maps
give elements of the respective second-coordinate lattices
on \(w^{-1}W_\tau\) and on \(H_\tau\subset w^{-1}T_\tau\).
Hence on these domains the graphs may be calculated with
\(i_*,q_*\) instead of the finite-level lifts, exactly.
Equation (AL3) now gives the first-coordinate image in
\(\Gamma_\tau\) equal to \(w^{-1}L_\tau\), and the second
coordinate image in \(\Delta_\tau\) equal to
\(O_\tau\). These are precisely (RC1), without an error
term or a uniform comparison of matrix entries.

For (RC2), when the chart graph is the actual first
differential graph, its pins in the preceding section
hold directly: products of two functions vanishing at
zero have zero differential, the character projector
commutes with differentiation, and
\(\omega_{P[p]}=\omega_P/p\omega_P\) is the quotient of
the augmentation ideal by its square. Apply the exact
cotangent-sequence and coordinate-ideal argument of that
section, using Corollary~\ref{cor:geometry-induced-additive}
for quotient generation.

This supplies the six test identities for the canonical
induced orders and differential graphs at the required
height-one fields. To use them with the previous global
relative-Chern comparison, one must identify the pullbacks
of the original chart function and differential lattices
with these orders and graphs (including at new height-one
valuations after flattening). Equivalently one can construct
an alternative coherent diagram on the same generic flag
variety with these test properties, keeping the original
four root bundles and their metrics. The original
\texttt{testOrd}, \texttt{testOrdW}, \texttt{funLat}
and \texttt{chartD} identifications are the outstanding
global interface; an equality of abstract finite-group
orders or of their ranks alone is not that identification.
No assertion of the weighted tower estimate or the
completion of W10 is made in this local argument.
